\section*{Chapitre \ref{ch:m1:options-markets-europe-asia} --- Marchés d'options en Europe et en Asie}

\begin{solution}{exo:m1:options-markets-europe-asia:1}
Notionnel $25 \times 22\,000 \times 0.012 = \$6\,600$~; prime $25 \times 35 \times
0.012 = \$10.50$.
\end{solution}

\begin{solution}{exo:m1:options-markets-europe-asia:2}
Notionnel 560\,000~\$~; prime 1\,600~\$. Un tel contrat porte une prime
équivalant à 152 petits contrats, et un notionnel équivalant à 85.
\end{solution}

\begin{solution}{exo:m1:options-markets-europe-asia:3}
30 milliards. Avec 24 milliards observés, l'activité en notionnel est $24 \times 75/(90
\times 25) = 80\%$ de ce qu'elle était~: un cinquième de l'activité
est passé aux intervenants plus modestes.
\end{solution}

\begin{solution}{exo:m1:options-markets-europe-asia:4}
2024~: $170.2 + 19.8 + 10.0 + 4.2 + 2.3 = 206.4$ milliards~; Asie-Pacifique 82~\%.
2025~: $75.59/119.29 = 63\%$. Vérification~: $1 - 119.29/206.4 = 42.2\%$.
\end{solution}

\begin{solution}{exo:m1:options-markets-europe-asia:5}
M1~: $945 \times 4\% \times \tfrac14 = \$9.45$ milliards. M2~: $1\,360 \times 1.5\%
\times \tfrac14 = \$5.1$ milliards. Le pool de M1 est plus grand malgré son
\omterm{def:m1:options-markets-europe-asia:turnover}{volume en primes} plus faible, parce que les fourchettes sont larges. Avant d'entrer~: frais et taxes~;
qui peut négocier et sous quelle forme juridique~; quelle part du pool les
acteurs établis mieux placés captent déjà~; la \omterm{def:m1:what-a-trading-firm-does:adverse-selection}{sélection adverse} dans chaque flux~; et
combien de temps la spécification du contrat restera telle quelle.
\end{solution}

\begin{solution}{exo:m1:options-markets-europe-asia:6}
Celui qui veut faire bouger le cours doit maintenir la pression pendant trente minutes au lieu d'un seul cours publié,
face à des arbitragistes qui ont le temps de peser en sens inverse~; et le cours de chaque
minute compte pour un trentième. Coût~: personne ne peut négocier \emph{au} prix de
règlement~: les hedgers qui veulent sortir au règlement doivent exécuter un
programme sur la fenêtre et supporter un \omterm{def:m1:the-buy-side:benchmark}{tracking error} par rapport à la moyenne.
\end{solution}

\begin{solution}{exo:m1:options-markets-europe-asia:7}
Avant~: contrats 98.1 / 0.9 / 0.4 / 0.5~\%~; notionnel 48.3 / 38.7 / 1.6 /
11.4~\%~; prime 38.2 / 54.9 / 6.9 / 0~\%. Après (M1 à 24 milliards de contrats de
\omterm{def:m1:options-markets-europe-asia:turnover}{taille de lot} 75)~: la part de M1 en contrats tombe à environ 93~\%, toujours écrasante,
tandis que son notionnel et sa prime, chacun à 80~\% de leur niveau antérieur, donnent des parts d'environ 43~\%
et 33~\%. Le tableau des contrats ne remarque presque pas un changement qui a retiré un cinquième de
l'activité du marché.
\end{solution}

\begin{solution}{exo:m1:options-markets-europe-asia:8}
La combinaison. Un régulateur, et un tribunal, examineront la société comme une seule
entité~: une vaste position vendeuse en options sur indice dont le gain dépend de
l'indice de clôture, détenue pendant que les transactions de la même société sur les composants
font bouger cet indice, d'abord à la hausse pendant la vente des options, puis à la baisse jusqu'au
règlement. La logique locale de chaque jambe ne répond pas à la question de
l'intention pour l'ensemble, et les propres registres de la société montrent cet ensemble. Ce qui
manque, c'est un contrôle qui agrège l'exposition au règlement à travers
les desks, limite les transactions sur le sous-jacent de celui qui profite de leur effet
sur un dérivé, et documente les raisons \emph{avant} la transaction.
\end{solution}

\begin{solution}{pb:m1:options-markets-europe-asia:1}
\textbf{1.} $90/0.85 = 106$ contre 1.
\textbf{2.} M1~: $90 \times 10^9 \times 6\,600 = 594$ billions de dollars. M2~: $0.85 \times
10^9 \times 560\,000 = 476$ billions de dollars.
\textbf{3.} M1~: $90 \times 10^9 \times 10.5 = \$945$ milliards. M2~: $0.85 \times 10^9
\times 1\,600 = \$1\,360$ milliards.
\textbf{4.} 1.44.
\textbf{5.} 15.9 et 28.6 points de base. Les options de M1 sont moins chères par rapport à
leur sous-jacent~: à échéance plus courte et plus loin hors de la monnaie, des
billets de loterie sur le mouvement du jour.
\textbf{6.} M1~: $945 \times 1\% = \$9.45$ milliards. M2~: $1\,360 \times 0.375\% =
\$5.10$ milliards.
\textbf{7.} Frais $0.05\% \times 945 = \$0.47$ milliard~; taxe sur la moitié de
la prime que vendent les \omterm{def:m1:what-a-trading-firm-does:market-maker}{teneurs de marché}, $0.1\% \times 472.5 = \$0.47$ milliard~: 0.94~milliard de dollars.
\textbf{8.} $0.85 \times 10^9 \times 0.50 = \$0.42$ milliard.
\textbf{9.} 8.51~milliards de dollars et 4.67~milliards de dollars.
\textbf{10.} Modeste au départ. Cela dépend de l'accès (adhésion,
co-localisation, statut de \omterm{def:m1:what-a-trading-firm-does:market-maker}{teneur de marché} et ses privilèges), de la rapidité et de la qualité
des prix par rapport aux acteurs établis, et de la part du flux qu'il voit~:
la fourchette moyenne du pool se gagne sur le flux bénin que les acteurs
titulaires de droits voient en premier.
\textbf{11.} $24 \times 75/(90 \times 25) = 80\%$.
\textbf{12.} $24 \times 10^9 \times 75 \times 35 \times 0.012 = \$756$ milliards.
\textbf{13.} $756 \times (1\% - 0.1\%) = \$6.8$ milliards, soit un cinquième de moins~; et
les participants qui sont partis étaient les moins informés, si bien que les
fourchettes gagnées sur le reste risquent d'être plus faibles aussi.
\textbf{14.} Une activité bâtie sur le flux de détail d'un seul marché porte un risque
réglementaire que nul modèle ne couvre~: le régulateur peut modifier la \omterm{def:m1:options-markets-europe-asia:turnover}{taille de lot},
les échéances, les marges ou l'accès en un trimestre. Dimensionner les coûts fixes
au marché qui survivrait à un tel changement, et détenir un second marché.
\textbf{15.} Son étude constate que 93~\% des traders particuliers perdent de l'argent, au
total plus de 1.8~billion de roupies sur trois ans, dont une grande part sur des options
le jour d'échéance. Des contrats plus grands, moins d'\omterm{def:m1:options-markets-europe-asia:weekly}{échéances hebdomadaires}, une prime
payée d'avance et une marge supplémentaire le jour d'échéance font tous monter le prix du
pari le moins cher.
\textbf{16.} D'abord les coûts de transaction, les courtiers et les bourses~; puis les
contreparties mieux loties en tarification et en exécution~: les \omterm{def:m1:what-a-trading-firm-does:prop}{sociétés de trading
pour compte propre} et les investisseurs étrangers, en grande partie algorithmiques.
\textbf{17.} Non~: coter des cours acheteur et vendeur à des acheteurs consentants est le service.
La limite est franchie lorsqu'un participant négocie le sous-jacent afin de
faire bouger le règlement dont dépendent ses options~; la proposition montre que
la tentation est structurelle partout où le marché des options éclipse son
sous-jacent, et c'est pourquoi la surveillance l'est aussi.
\textbf{18.} Une voie juridique (catégorie d'enregistrement, entité locale ou
partenaire), un membre compensateur, une structuration fiscale, des \omterm{def:m1:options-markets-europe-asia:weekly}{limites de position}
adaptées à la stratégie, une équipe locale de conformité et une compréhension de
la façon dont le régulateur regarde le trading pour compte propre étranger.
\textbf{19.} \textbf{1.44~: le marché qui détient 1~\% des contrats a un \omterm{def:m1:options-markets-europe-asia:turnover}{volume en primes}
supérieur de 44~\%.}
\textbf{20.} Le nombre de fois où un moteur d'appariement a publié une transaction, ce qui
constitue la source de revenus de la bourse et de personne d'autre.
\end{solution}

\begin{solution}{iq:m1:options-markets-europe-asia:1}
Selon quelle mesure~: nombre de contrats, notionnel ou primes~? Quelle est la \omterm{def:m1:options-markets-europe-asia:turnover}{taille de lot} et
comment a-t-elle évolué~? Quelle part représentent les options face aux contrats à terme, et les indices face aux
actions isolées, et quelle part expire le jour même~? Qui négocie~: la part des particuliers
et celle des sociétés pour compte propre~? Quel serait le classement
en volume de primes, c'est-à-dire en argent~.
\iqlookfor{les trois mesures, sans indication.}
\end{solution}

\begin{solution}{iq:m1:options-markets-europe-asia:2}
Une seule bourse et une seule chambre de compensation par produit, et non plusieurs bourses sur une
même chambre de compensation~: pas de fongibilité entre plateformes, pas de \omterm{def:m1:us-equity-market-structure:nbbo}{meilleurs cours acheteur et vendeur nationaux} (NBBO),
pas d'\omterm{def:m1:options-market-structure:auction}{enchères d'amélioration de prix}. Les options sur indices sont de style européen et \omterm{def:m1:basis-roll-and-delivery:delivery}{réglées en
espèces}, couvertes par le contrat à terme sur l'indice local. Une grande part de la taille institutionnelle
se traite sous forme de blocs négociés. La demande des particuliers pour l'effet de levier se tourne surtout vers
les warrants et certificats émis par les banques, et non vers les options cotées.
\iqlookfor{la structure de compensation comme cause première de la différence.}
\end{solution}

\begin{solution}{iq:m1:options-markets-europe-asia:3}
Le nombre de contrats baisse mécaniquement d'au moins un facteur trois. Le notionnel et
les primes baissent moins, à hauteur du nombre de participants évincés~: les plus petits
comptes, dont le ticket minimum a triplé. Le flux restant est un peu
mieux capitalisé et mieux informé~; les spreads exprimés en primes pourraient
se resserrer. Les \omterm{def:m1:what-a-trading-firm-does:market-maker}{teneurs de marché} perdent le flux le plus bénin~; les bourses perdent les recettes de frais
par unité de notionnel si les frais sont prélevés par contrat. Une partie de l'activité
migre vers d'autres produits à effet de levier.
\iqlookfor{effets mécaniques face aux effets comportementaux.}
\end{solution}

\begin{solution}{iq:m1:options-markets-europe-asia:4}
Leur gain dépend d'un chiffre calculé à partir de quelques cours sur une courte
fenêtre, dans un marché bien plus petit que les positions dérivées qui le prennent pour
référence, si bien qu'un léger mouvement du sous-jacent peut valoir plus que le
coût engagé. Les parades sont un règlement sur une moyenne calculée sur une fenêtre, ou sur une enchère
avec informations de déséquilibre et bornes de prix~; des \omterm{def:m1:options-markets-europe-asia:weekly}{limites de position} suivies
en cours de séance le jour d'échéance~; des marges plus élevées ce jour-là~; une surveillance qui
rapproche les activités au comptant et sur dérivés des participants~; et étale les échéances
loin des sous-jacents isolés et peu liquides.
\iqlookfor{le décalage de taille et au moins trois parades.}
\end{solution}

\begin{solution}{iq:m1:options-markets-europe-asia:5}
Une vue consolidée, en temps réel, de l'exposition à la valeur de règlement, sur
l'ensemble des desks et des entités. Une règle selon laquelle celui qui profite de la variation
de l'indice ne peut pas négocier agressivement les composants pendant la fenêtre de règlement,
ou seulement selon un calendrier préapprouvé. Une justification écrite, préalable à toute grosse
transaction au comptant le jour d'échéance. Des limites sur la part de la société dans le volume des
composants pendant la fenêtre. Une revue indépendante de l'\omterm{def:m1:pnl-and-positions:attribution}{attribution du P\&L}
le jour d'échéance. Et une escalade lorsqu'une bourse ou un régulateur pose des questions.
\iqlookfor{agrégation entre desks, et documentation avant la transaction.}
\end{solution}

\begin{solution}{iq:m1:options-markets-europe-asia:6}
Taille en argent~: \omterm{def:m1:options-markets-europe-asia:turnover}{volume en primes}, spreads, donc le réservoir de revenus~; puis les frais
et les taxes. Flux~: qui négocie, avec quelle information, à quel point il se concentre dans le temps
(jours d'échéance). Accès~: forme juridique, adhésion, colocalisation, programmes de \omterm{def:m1:what-a-trading-firm-does:market-maker}{teneurs de marché}
et leurs obligations, \omterm{def:m1:options-markets-europe-asia:weekly}{limites de position}. Concurrence~: part des acteurs en place
et avantage concurrentiel. Politique~: préoccupations affichées du régulateur, changements récents et en cours
sur les \omterm{def:m1:options-markets-europe-asia:turnover}{tailles de lot}, les échéances et les marges, attitude envers les sociétés pour compte
propre étrangères~; scénario avec un réservoir réduit de moitié. Opérations~: compensation, collatéral,
devise, fiscalité, personnel. Sortie~: coûts fixes récupérables si les règles changent.
\iqlookfor{risque politique traité comme un élément de premier plan, avec un scénario.}
\end{solution}
